\chapter{来源与史料说明}
\label{app:sources}

本附录说明本卷索引所指向的材料路径，而不是把文献排成一张不变的高低表。隋唐五代的官方纪事、制度文本、地方与出土材料保存状况差异很大：它们各能回答不同问题，也各会放大不同的政治视角。阅读事件时，应先辨认材料的体裁、成书与保存条件，再判断它能够确认的行动和必须保留的空白。事件入口见\hyperref[app:index]{“隋唐五代索引”}。

\section{正史与编年的年代脊柱}
\label{appsrc:chronology}

\begin{description}
  \item[隋代] 《隋书》帝纪、志与列传，《北史》以及《资治通鉴》共同提供581至618年的基本年代框架。帝纪善于追踪诏令、战争和宫廷政变；志书与列传则可补官制、礼制、地理和人物网络。它们不能单独重建地方社会的全部负担，尤其不能把后出的整齐叙事当作现场记录。
  \item[唐代] 《旧唐书》《新唐书》与《资治通鉴》是追踪建国、武周、开元、安史与晚唐政治线的核心入口。三者的取材、编纂时间和叙述目标不同；对同一人物的褒贬、对藩镇和宦官的解释、以及战时消息的先后，都应相互核对。
  \item[五代十国] 《旧五代史》《新五代史》与《资治通鉴》保留中原政权更替的主线，但不能让中原朝廷的年号和正统观覆盖南方诸国、边境政权或地方社会。国别记载、墓志与区域文书须与之并读。
\end{description}

\section{制度、法律与行政材料}
\label{appsrc:institutions}

《唐六典》《唐会要》《通典》《唐律疏议》以及诏敕、制书、表奏和相关类书，是理解官制、礼仪、财政、法律和行政语言的重要入口。《册府元龟》可以帮助追索后世汇编的材料线索，但其编排不等于原始档案的时间顺序。制度文本首先说明规范如何被表述；是否、何时和在何地实行，仍须同地方文书、墓志、账簿材料与具体事件并读。

\section{边疆、宗教与跨区域交流}
\label{appsrc:frontiers}

敦煌、吐鲁番等地文书，墓志、碑刻、造像题记与钱币，可使官修史书之外的军镇、寺院、商旅、家庭和书吏进入视野。对突厥、吐蕃、高句丽、新罗、日本及中亚的行动，还应同相应区域的传世与出土材料互证；译名、纪年和政治称号并不总能一一对应。这些材料尤其适合检验人员流动、交通节点和文书实践，却不能由个别遗址或一件文书推断整个帝国的统一经验。

\section{城市、运河与物质材料}
\label{appsrc:material}

大兴—长安、洛阳、运河沿线城址与水工遗存，以及墓葬、窖藏和交通遗迹，是理解都城布局、运输条件、手工业与区域联系的关键证据。考古发现可以校正文本中的空间想象和工程规模，却通常难以直接说明某项政策的决策过程、劳动者的全部经历或一场战争中每支军队的行动。使用时应留意发掘背景、断代和发表材料的限度。

\section{五代十国的区域材料}
\label{appsrc:five-dynasties}

五代十国不宜只按“中原朝代之间的间隙”阅读。除两《五代史》和《资治通鉴》外，国别史料、地方志的早期材料层、墓志、寺院碑刻、钱币与城市考古，都有助于辨认吴越、南唐、荆南、闽、南汉、前后蜀及北方边境政权的不同资源与记忆。它们的保存往往零散，恰须避免把后来的区域认同、地名或行政边界倒投回十世纪。

\section{怎样沿索引复核}

可先在\hyperref[app:index]{“隋唐五代索引”}按事件、来源或图件入口进入本卷，再回到本附录判断材料适用范围。一个事件有多个来源，并不自动提高每个细节的确定性：重要的是确认它们是否独立、是否接近所问的问题，以及分歧会不会改变对行动、制度或社会代价的理解。
